\section*{Chapter \ref{ch:rs:alternative-and-text-data} --- Alternative and Text Data}

\begin{solution}{exo:rs:alternative-and-text-data:1}
$1\,200/3\,000 = 40\%$ of the rows describe periods before the launch and may have been built afterwards. Ask for the first delivery date of
every row (or the archived files as delivered), the launch date, how the pre-launch history was constructed, and every change of source or
method with its date.
\end{solution}

\begin{solution}{exo:rs:alternative-and-text-data:2}
General list: $12/400 = 3\%$. Without the four finance words: $3/400 = 0.75\%$. Net tone with the second list: $(3 - 5)/400 = -0.5\%$.
\end{solution}

\begin{solution}{exo:rs:alternative-and-text-data:3}
$4 \times 2.9\% \times \$20$ million $= \$2.32$ million, less $0.4\% \times \$20$ million $= \$0.08$ million: \$2.24 million a year.
\end{solution}

\begin{solution}{exo:rs:alternative-and-text-data:4}
Both measure the same $s$, so they are correlated; residualising $m$ on $o$ removes the part of $m$ predictable from $o$, including part of
its $s$ component, and what remains is the part of $s$ that $o$ misses plus $m$'s own noise: positive, smaller. If $e_2$ were a copy of $e_1$,
$m$ would be a rescaled copy of $o$ (plus nothing), its residual would be zero, and the incremental IC nil.
\end{solution}

\begin{solution}{exo:rs:alternative-and-text-data:5}
The profit of year $k$ (from 0), measured at its middle, is $A\,2^{-(k+1/2)/h}$; the fee equal to the mean annual profit is the mean over the
$Y$ years. For $A = \$2.3$ million, $Y = 3$, $h = 2$: $\frac{2.3}{3}(0.841 + 0.595 + 0.420) = \$1.42$ million.
\end{solution}

\begin{solution}{exo:rs:alternative-and-text-data:6}
The model's parameters were fitted on text written up to 2025, including the commentary, outcomes and hindsight about 2015 to 2020: which
companies failed, what a phrase turned out to mean. Its sense of which documents are alike, and its vocabulary, encode that knowledge. A
point-in-time backtest uses a model trained only on text available before each decision date, or at least measures the difference.
\end{solution}

\begin{solution}{exo:rs:alternative-and-text-data:7}
\texttt{rs\_altdata.ic\_break}: the quarterly ICs are 0.32 to 0.45 in the first eight quarters and $-0.03$ to 0.27 after; the CUSUM puts the break
at the ninth quarter, the start of the third year, with a statistic of 1.85, significant but far weaker than the 3.05 of the fit to the truth.
Twenty noisy quarterly ICs carry less information about the break than 5\,980 rows of the panel against the reported values: the fit measures the
panel's relation to the quantity it claims to measure directly, the IC only through the returns' noise.
\end{solution}

\begin{solution}{exo:rs:alternative-and-text-data:8}
Anonymisation is a representation, and App Annie made such representations while using non-anonymised data; the buyer needs to know how the data
were collected and under which consents, and to document its own diligence. Publicly accessible is not the same as free to use: a court's view of one
statute on access to public profiles does not settle contract terms, copyright, data-protection law (personal data under the GDPR include online
identifiers and location data) or whether any source was under a duty of confidence. Compliance signs off before the trial.
\end{solution}

\begin{solution}{pb:rs:alternative-and-text-data:1}
\begin{enumerate}
\item 5\,980 rows. A panel value delivered after the announcement cannot predict it: those rows are dropped rather than scored.
\item 338 a quarter in the backfilled years, 273 in the live ones. A panel whose history covers more names than its live product was built after
the fact on a different basis.
\item 39.6\%.
\item The first delivery date of each row, or the launch date and the statement of which history was reconstructed.
\item Slopes 0.92, 0.87, 0.37, 0.39, 0.35; correlations 0.85, 0.84, 0.35, 0.35, 0.32; intercepts 0.02 in the first two years, 0.12 to 0.17 after.
\item At the start of the third year, with a statistic of 3.05.
\item The vendor went live when it changed provider and reconstructed the previous years from the new provider's panel re-weighted to reported
results; the level, the noise and the coverage changed at the same date.
\item That the history before the break is of a different kind from the live data, whatever its cause, and that only the live years price the
product.
\item 0.38, 0.37, 0.12, 0.14, 0.15; the firm's own signal 0.20 in the live years.
\item 0.10, with a $t$ statistic of 4.0 over 12 quarters.
\item 10.4\% a quarter (backfilled), 3.4\% (live, raw), 2.9\% (live, incremental).
\item The backfill is fitted to outcomes and the firm will not receive data like it in future; what the firm already owns it does not pay for twice.
\item \$5.06 million, \$1.59 million and \$1.35 million a year.
\item \$0.43 million at six months, \$1.81 million at five years.
\item \textbf{Named result.} On live data, net of the firm's own signal, with \$20 million of capacity, 0.4\% a year of costs and a two-year
half-life, the card panel is worth \$1.35 million a year over three years, a quarter of what its backfilled history implies (\$5.06 million); at
a six-month half-life, \$0.43 million.
\item At \$400\,000, yes unless the edge is expected to halve within about six months; at \$900\,000, only if its half-life is at least about a year
(break-even \$0.87 million at one year): the decision is a bet on the number of other buyers.
\item The first-delivery archive and change log, a longer trial on live data, exclusivity or a cap on the number of buyers, notification of method
changes, and a price step-down if coverage falls.
\item How the data are collected and under which consents; whether any source owes a duty of confidentiality; whether personal data are present
and how they are protected; the sites' terms if scraped; the vendor's representations in writing; compliance sign-off.
\item Measure the live IC and the long--short spread quarter by quarter from purchase, compare them with the trial's, and estimate the decay
(chapters 13 and 28).
\item What it adds to what the firm already has, on data it will actually receive, net of costs, for as long as the edge lasts.
\end{enumerate}
\end{solution}

\begin{solution}{iq:rs:alternative-and-text-data:1}
Backfilled history built after the fact (and fitted to outcomes); \omterm{def:rs:alternative-and-text-data:altdata}{panel drift} (provider or method changes); timestamps that are periods, not
deliveries; entity mapping with today's identifiers; revisions without vintages; survivorship in the covered names.
\end{solution}

\begin{solution}{iq:rs:alternative-and-text-data:2}
Residualise the new feature on the existing ones in each cross-section and measure the IC of the residual (the incremental IC) and its $t$
statistic, on live data only; better, add it to the production combination (chapter~14) and measure the change in the combined forecast's
performance out of sample.
\end{solution}

\begin{solution}{iq:rs:alternative-and-text-data:3}
The document's publication time (not the event it describes); the model's training data must predate each decision date, or the model knows the
outcome; the model version and its tokenizer change the scores, so they are stored with the scores; entity linking with point-in-time
identifiers; revisions and deletions of articles.
\end{solution}

\begin{solution}{iq:rs:alternative-and-text-data:4}
Estimate the incremental profit: live data only, the incremental feature's long--short return, the capacity the firm can deploy, costs, and a decay
rate that reflects how many others buy; the break-even fee is the mean net annual profit over the contract. Pay if \$500\,000 is comfortably below it
under a pessimistic half-life, after legal clearance.
\end{solution}

\begin{solution}{iq:rs:alternative-and-text-data:5}
It may contain \omterm{def:rs:alternative-and-text-data:mnpi}{material non-public information} obtained in breach of a duty, personal data, or data collected against terms or consents; the SEC has
charged a data vendor for misrepresenting how its data were derived. Diligence: written provenance and consents, confidentiality of sources,
personal-data review, terms of use, compliance approval, and records of all of it.
\end{solution}

\begin{solution}{iq:rs:alternative-and-text-data:6}
Their lists count words that are negative in general usage but neutral in finance (Loughran and McDonald found almost three-fourths of the Harvard
dictionary's negative words in 10-Ks were such words). Use finance-specific lists, net of positive words, relative to the firm's own past documents;
or a model trained on financial text available at the decision date, validated against returns or a labelled sample.
\end{solution}
